\section{Candidate follow-up: EXP-009}
\label{sec:followup}

We sketch one natural candidate experiment that, if eventually run,
would partition the two hypotheses behind the DermLIP positive. We
present the design now to invite feedback on whether this is the right
follow-up before we (or anyone else) invests the engineering effort
required to run it.

\subsection{The partition question}
The DermLIP positive in EXP-007 ($\auroc = 0.944 \pm 0.003$) is consistent
with two distinct hypotheses. Either DermLIP's dermoscopy-domain pretraining
on Derm1M's non-HAM10000 fraction has aligned nuisance directions across the
three families and that alignment generalises to \strongholdtwo, or DermLIP
saw HAM10000 raw images during pretraining and the embedding space is
well-fit to HAM10000 lesion identity in a way that survives our synthetic
nuisance. EXP-008's BiomedCLIP partition rules out only the broader ``any
medical-image pretraining'' alternative. The dermoscopy-vs-HAM10000
distinction remains open.

\subsection{EXP-009 design}
Self-pretrain a DINOv2 ViT-B/14 encoder on a non-HAM10000 dermoscopy
corpus assembled from ISIC archives minus the HAM10000 split, DermNet,
DermQuest, BCN20000 \citep{combalia2019bcn20000} non-HAM10000 components, and
similar publicly available dermoscopy and clinical-skin sources. Run a short
JEPA-style or masked-image-modelling objective (target: $\sim$20 epochs,
$\sim$200\,k steps), then re-run the EXP-004 recipe on top with the linear
predictor.

\subsection{Decision table}
Table~\ref{tab:exp009-decision} maps the EXP-009 test-$\auroc$ band to a
revision of the paper's claim.

\begin{table}[h]
\centering
\small
\setlength{\tabcolsep}{4pt}
\begin{tabular}{@{}lp{0.32\linewidth}p{0.40\linewidth}@{}}
\toprule
EXP-009 test $\auroc$ & Interpretation & Effect on the paper claim \\
\midrule
$\geq 0.85$    & Dermoscopy-domain transfer is sufficient. & Headline survives; caveat downgrades to ``verified.'' \\
$0.50$--$0.80$ & Partial transfer; mixed contribution. & Claim becomes ``domain helps, dataset overlap helps more.'' \\
$0.30$--$0.50$ & HAM10000 overlap drove EXP-007. & Claim narrows to ``frozen-backbone JEPA unlocks when the backbone has seen the eval dataset.'' \\
\bottomrule
\end{tabular}
\caption{Three EXP-009 outcome bands on \strongholdtwo\ and the
corresponding revision of the paper's claim. Committing this mapping up
front would prevent post-hoc re-interpretation of the experiment.}
\label{tab:exp009-decision}
\end{table}

\subsection{Why we share the design before running it}
Corpus assembly is multi-week: provenance audit on every ISIC component
to confirm HAM10000 exclusion, SSL pretrain compute on a non-trivial
GPU budget, and a pretrain-quality check on a held-out HAM10000 split.
Releasing the methodology and the nine-run arc at this stage is
intentional --- the paper seeks community feedback on the proxy-task
design and on whether this partition is the right follow-up before
anyone (including us) commits to the build. Whether EXP-009, a
different partition, or no follow-up at all is the right next step
depends on that feedback.

\subsection{Feedback we are seeking}
\begin{enumerate}[leftmargin=2em]
\item Is the proxy-task design (three deterministic nuisance families with
      operation-level disjointness, train-on-two / evaluate-on-third) a
      reasonable test of frozen-backbone generalisation, or is a different
      construction more informative?
\item Is the proposed partition (a non-HAM10000 dermoscopy SSL pretrain) the
      right one, or would a different partition (e.g.\ a dermoscopy-text-only
      retrieval probe over DermLIP, or a fine-tuning-on-HAM10000 ablation of
      a known-non-contaminated backbone) be more informative?
\item What other follow-ups would matter for the next stage --- different
      architecture classes, real same-lesion-over-time data, or a different
      evaluation metric?
\item Are there pre-trained backbones in the public literature that would be
      cleaner partition points than the ones we have used?
\end{enumerate}
